% Chapter 3: Arbitrage in the strike and the smile.

Fix a maturity and look at the price of a European call as a function of its strike. Before
any model is chosen, absence of arbitrage already shapes this function: it is convex and
nonincreasing, its slopes lie between $-e^{-rT}$ and $0$, and it starts at the price of the
stock without its dividends. Minus its first derivative is the price of a digital option, and its
second derivative is the discounted risk-neutral density of $S_T$: the
\emph{Breeden--Litzenberger formula}, which turns a strip of option prices into a probability
law. Chapter~1 met the same object as the Green's function of the pricing equation
(\cref{prop:sm-green}).

Markets quote options through their Black--Scholes implied volatilities, and for equity indices
these are far from flat: they fall as the strike rises. Since the strike enters a price through
the implied volatility as well, every strike derivative picks up terms in the slope and the
curvature of the smile. For an at-the-money digital the first of them, the \emph{skew term}, is
about a fifth of the price (\cref{exa:sk-course-digital}).

The IEOR 222 notes prove the strike relations in a one-period market without dividends. Here we
add the yield, prove the Breeden--Litzenberger formula in the generality that market data need
(atoms in the law of $S_T$, call curves with kinks), characterize exactly which functions are call
prices, and turn to methods: extracting a density from quotes and choosing the strike spacing
(\cref{prop:sk-fd-density}), pricing and replicating digitals under a smile
(\cref{prop:sk-binary-smile,prop:sk-call-spreads}), and testing a parametric smile for
arbitrage. \Cref{ch:04-static-replication} builds general payoffs out of the same strip of
options.

Throughout, $T>0$ is a fixed maturity, $r,y\in\R$, $S_0>0$, $F=S_0e^{(r-y)T}$ is the forward
price, and $\Delta_h^2f(K)=\big(f(K-h)-2f(K)+f(K+h)\big)/h^2$ is the second difference of a
function $f$ with step $h>0$. One-sided derivatives are written $\partial^\pm f(K)$.

\section{Call prices across strikes}\label{sec:sk-strike}

\paragraph{The market.} At time $0$ the following claims on the date $T$ are traded, in a
frictionless market (any real quantity, no bid--ask spread):
\begin{itemize}
  \item the zero-coupon bond, which pays $1$ and costs $e^{-rT}$;
  \item the \emph{prepaid forward}, which pays $S_T$ and costs $S_0e^{-yT}$; in the standard model
    of \cref{def:sm-model}, $e^{-yT}$ shares bought at time $0$, with the dividends reinvested in
    shares, grow to one share at $T$;
  \item European calls and puts with strikes $K\ge0$, which pay $(S_T-K)^+$ and $(K-S_T)^+$ and
    cost $C(K)$ and $P(K)$.
\end{itemize}
The terminal price $S_T$ is a random variable on $(\Omega,\F,\Prob)$ with $S_T\ge0$ (limited
liability). A static portfolio holds fixed quantities of finitely many of these claims; its
price and its payoff are the corresponding linear combinations. We use one of two hypotheses.
\begin{enumerate}[label=(\Alph*)]
  \item \emph{No static arbitrage:} no static portfolio with a payoff $V\ge0$ almost surely has
    a negative price, or a price of zero together with $\Prob(V>0)>0$.
  \item \emph{Pricing measure:} there is a probability measure $\Q$ equivalent to $\Prob$ with
    $\EQ[S_T]<\infty$ under which every traded payoff $X$ costs $e^{-rT}\EQ[X]$. In particular
    $\EQ[S_T]=e^{rT}S_0e^{-yT}=F$, and $C(K)=e^{-rT}\EQ[(S_T-K)^+]$.
\end{enumerate}
Both imply \emph{positivity}:
\begin{enumerate}[label=(P)]
  \item a static portfolio whose payoff is nonnegative almost surely has a nonnegative price.
\end{enumerate}
Under (A) this is part of the definition; under (B) it holds because the expectation of a
nonnegative variable is nonnegative, and $\Prob$ and $\Q$ have the same null sets. Applied to a
payoff $V$ and to $-V$, (P) gives the \emph{law of one price}: if two portfolios have the same
payoff almost surely, they have the same price. Hypothesis (B) implies (A) (the argument is
the one for a one-period market in the IEOR 222 notes); the converse, the fundamental theorem of
asset pricing, is not needed here. The proposition below uses (P) only, so it is
\emph{model-free}: it holds in every market free of static arbitrage, whatever the law of $S_T$.

\begin{proposition}[Strike relations]\label{prop:sk-strike-bounds}
Assume (P). Let $K\ge0$, $0\le K_1<K_2$, $\lambda\in[0,1]$ and $K_\lambda=\lambda K_1+(1-\lambda)K_2$.
\begin{enumerate}[label=(\roman*)]
  \item (\emph{Parity and bounds}) $C(K)-P(K)=S_0e^{-yT}-Ke^{-rT}$, $C(0)=S_0e^{-yT}$, $P(0)=0$, and
    \[
      \big(S_0e^{-yT}-Ke^{-rT}\big)^+\le C(K)\le S_0e^{-yT},\qquad
      \big(Ke^{-rT}-S_0e^{-yT}\big)^+\le P(K)\le Ke^{-rT}.
    \]
  \item (\emph{Monotonicity and slopes})
    $-e^{-rT}(K_2-K_1)\le C(K_2)-C(K_1)\le0$ and $0\le P(K_2)-P(K_1)\le e^{-rT}(K_2-K_1)$.
  \item (\emph{Convexity}) $C(K_\lambda)\le\lambda C(K_1)+(1-\lambda)C(K_2)$, and the same for $P$.
\end{enumerate}
Consequently $C$ is convex, nonincreasing and Lipschitz with constant $e^{-rT}$ on $[0,\infty)$.
It has one-sided derivatives at every $K>0$, with
$-e^{-rT}\le\partial^-C(K)\le\partial^+C(K)\le0$, and a right derivative $\partial^+C(0)\in[-e^{-rT},0]$.
The put is convex and nondecreasing, with slopes in $[0,e^{-rT}]$, and $P''=C''$ wherever either exists.
\end{proposition}

\begin{proof}
For $x\ge0$ and $K,K_1,K_2$ as in the statement, the following hold:
\begin{gather}
  (x-K)^+-(K-x)^+=x-K,\qquad 0\le(x-K)^+,\qquad x-K\le(x-K)^+\le x,\label{eq:sk-payoff-bounds}\\
  (x-K_1)^+-(x-K_2)^+=\min\big((x-K_1)^+,\,K_2-K_1\big)\in[0,K_2-K_1],\label{eq:sk-call-spread}\\
  \lambda(x-K_1)^++(1-\lambda)(x-K_2)^+-(x-K_\lambda)^+\ge0 .\label{eq:sk-fly}
\end{gather}
The upper bound in \eqref{eq:sk-payoff-bounds} uses $x\ge0$ and $K\ge0$; \eqref{eq:sk-call-spread}
is checked on $x\le K_1$, $K_1<x<K_2$ and $x\ge K_2$; and \eqref{eq:sk-fly} holds because
$K\mapsto(x-K)^+$ is the maximum of the affine functions $0$ and $x-K$, hence convex. At $x=S_T$
each line compares the payoffs of two static portfolios of calls, puts, bonds and the prepaid
forward, and (P) with the law of one price turns it into the claimed relation between prices.
For example, a call minus a put pays $S_T-K$, which is also the payoff of the prepaid forward
minus $K$ bonds; this is parity. A call with strike $0$ pays $S_T$, so $C(0)=S_0e^{-yT}$, and a put
with strike $0$ pays nothing. The \emph{call spread} \eqref{eq:sk-call-spread} pays between $0$
and $K_2-K_1$, the payoff of $K_2-K_1$ bonds; this is (ii) for calls. The portfolio in
\eqref{eq:sk-fly} is a \emph{butterfly}; this is (iii). By parity $P$ differs from $C$ by the
affine function $K\mapsto Ke^{-rT}-S_0e^{-yT}$, which turns the relations for calls into those for
puts and leaves second derivatives unchanged.

For the derivatives, (iii) says that the slope $(C(b)-C(a))/(b-a)$ of a chord is nondecreasing in
each endpoint, for $0\le a<b$. So as $h\downarrow0$ the quotients $(C(K+h)-C(K))/h$ decrease and
$(C(K)-C(K-h))/h$ increase, the first stays above the second, and by (ii) both lie in $[-e^{-rT},0]$.
Their limits are $\partial^+C(K)\ge\partial^-C(K)$ (only the first one at $K=0$).
\end{proof}

Read backwards, each line of the proof is a trade. A call spread quoted at a negative price, or
above $e^{-rT}(K_2-K_1)$, is bought, respectively sold against $K_2-K_1$ bonds; a butterfly quoted
at a negative price is bought. The IEOR 222 notes work these trades out in detail.

\begin{remark}[Edges and atoms]\label{rem:sk-edges}
At $K=0$ the call is worth $S_0e^{-yT}$, not $S_0$: it is a claim to the share at $T$ without the
dividends paid before $T$. Under (B), the right slope at the origin is
$\partial^+C(0)=-e^{-rT}\Q(S_T>0)$ (\cref{thm:sk-breeden-litzenberger}(i)): it equals the extreme value
$-e^{-rT}$ exactly when $\Q(S_T=0)=0$, and a stock that can default with positive risk-neutral
probability has a flatter call curve at the origin (\exref{ex:sk-default}). As $K\to\infty$,
\cref{prop:sk-strike-bounds} only gives that $C$ decreases to some limit $\ge0$; under (B) the limit
is $0$, by dominated convergence with $(S_T-K)^+\le S_T$, and the slopes
$\partial^+C(K)=-e^{-rT}\Q(S_T>K)$ tend to $0$ as well. A kink of $C$ at some $K>0$, that is
$\partial^-C(K)<\partial^+C(K)$, is the signature of an atom of the law of $S_T$ at $K$ (again
\cref{thm:sk-breeden-litzenberger}(i)), such as a cash takeover offer at the price $K$.
\end{remark}

\begin{remark}[Finitely many strikes, several maturities]\label{rem:sk-finite}
Markets quote calls at finitely many strikes $0<K_1<\dots<K_n$. With $K_0=0$ and $C(K_0)=S_0e^{-yT}$,
let $s_i=(C(K_i)-C(K_{i-1}))/(K_i-K_{i-1})$. \Cref{prop:sk-strike-bounds} requires
\begin{equation}\label{eq:sk-checks}
  -e^{-rT}\le s_1\le s_2\le\dots\le s_n\le0\qquad\text{and}\qquad C(K_n)\ge0,
\end{equation}
and conversely \eqref{eq:sk-checks} implies every relation of the proposition between quoted strikes,
because the piecewise-linear interpolation of the quotes is then convex, nonincreasing and has
slopes in $[-e^{-rT},0]$. With equal spacing $h$, $s_{i+1}-s_i=h\,\Delta_h^2C(K_i)$ is the price of a
butterfly divided by $h$. These conditions are \emph{necessary}. Whether quotes satisfying them can
be produced by some law of $S_T$ with mean $F$ (a ``model'' in the sense of \textcite{davis-hobson},
which chooses the null sets too) is a separate theorem: yes, except in the boundary case
$s_n=0<C(K_n)$, a ``weak arbitrage'' (their Theorem~3.1 for one maturity, Theorem~4.2 for several;
the IEOR 222 notes construct a discrete law). Under (B) with a given $\Prob$, the $\Q$-law of $S_T$ must
be equivalent to its $\Prob$-law, which can rule out more. If the $\Prob$-law has full support, so has the
$\Q$-law, and then all inequalities in \eqref{eq:sk-checks} are strict: $s_1>-e^{-rT}$ because
$\Q(S_T<K_1)>0$, $s_{i+1}>s_i$ because the butterfly on $K_{i-1},K_i,K_{i+1}$ pays off on
$(K_{i-1},K_{i+1})$, and $s_n<0<C(K_n)$ because $\Q(S_T>K_n)>0$. Conversely, if the $\Prob$-law has a
positive density on $(0,\infty)$ (lognormal, say), strict inequalities suffice. Pick
$\sigma_0=-e^{-rT}$, $\sigma_i\in(s_i,s_{i+1})$ for $i<n$ and $\sigma_n\in(s_n,0)$; the convex function through the
quotes with slopes $\sigma_0<\sigma_1<\dots<\sigma_n<0$, then $0$, has one kink in each $(K_{i-1},K_i)$ and one in
$(K_n,\infty)$, and by \cref{thm:sk-breeden-litzenberger}(iv) it prices calls for a law with an atom at each
kink. Replace each atom by a positive density on its cell with the same mass and mean (a small weight
on a fixed positive density on the cell, the rest uniform near the atom, will do): the prices at the $K_i$ do not change, since every $(x-K_i)^+$ is affine on
each cell. With $q$ the resulting density and $p$ the $\Prob$-density of $S_T$, $\dd\Q/\dd\Prob=q(S_T)/p(S_T)$
gives (B). For a rectangular grid (the same strikes at every maturity), with zero rates and
dividends, \textcite{carr-madan-2005} show that nonnegative prices of call spreads, butterflies and
calendar spreads suffice to exclude static arbitrage.
\end{remark}

\begin{remark}[Calendar spreads]\label{rem:sk-calendar}
Let calls with maturities $T_1<T_2$ be priced by one pricing measure under which
$X_t=e^{-(r-y)t}S_t$ is a martingale on $[0,T_2]$, as in the standard model. Then
$C(K,T)=e^{-yT}\EQ[(X_T-\kappa)^+]$ with $\kappa=Ke^{-(r-y)T}$, and conditional Jensen,
$\EQ[(X_{T_2}-\kappa)^+\mid\F_{T_1}]\ge(X_{T_1}-\kappa)^+$, gives
\begin{equation}\label{eq:sk-calendar}
  e^{yT_2}C\big(\kappa e^{(r-y)T_2},T_2\big)\ \ge\ e^{yT_1}C\big(\kappa e^{(r-y)T_1},T_1\big),\qquad\kappa\ge0 :
\end{equation}
calls are nondecreasing in the maturity when the strike moves with the forward and prices are
counted in prepaid forwards. In terms of implied volatilities (\cref{sec:sk-smile}), the total
implied variance $\sigma^2T$ at a fixed ratio $K/F$ is nondecreasing in $T$. At a fixed strike
the call can lose value with maturity when $y>0$ (\exref{ex:sk-calendar}).
\end{remark}

\section{The Breeden--Litzenberger formula}\label{sec:sk-bl}

In this section (B) holds, and $\mu$ is the law of $S_T$ under $\Q$, a probability measure on
$[0,\infty)$ with mean $F$. It is convenient to define $C(K)=e^{-rT}\EQ[(S_T-K)^+]$ for every real
$K$: for $K<0$ this is $S_0e^{-yT}-Ke^{-rT}$, the price of the prepaid forward and $-K$ bonds. The
butterfly with strikes $K-h$, $K$, $K+h$ pays, for every $x$,
\begin{equation}\label{eq:sk-tent}
  (x-K+h)^+-2(x-K)^++(x-K-h)^+=h^2\Lambda_h(x-K),\qquad \Lambda_h(u)=\frac{(h-|u|)^+}{h^2},
\end{equation}
a tent of height $1/h$ and area $1$ centred at $K$ (check the four intervals cut by $K-h$, $K$,
$K+h$). Its price is $h^2\Delta_h^2C(K)$.

\begin{theorem}[Breeden--Litzenberger]\label{thm:sk-breeden-litzenberger}
Assume (B), and let $\mu$ be the law of $S_T$ under $\Q$.
\begin{enumerate}[label=(\roman*)]
  \item (\emph{Digitals}) For every $K>0$,
    \[
      \partial^+C(K)=-e^{-rT}\Q(S_T>K),\qquad \partial^-C(K)=-e^{-rT}\Q(S_T\ge K),
    \]
    and $\partial^+C(0)=-e^{-rT}\Q(S_T>0)$. So $\partial^+C(K)-\partial^-C(K)=e^{-rT}\Q(S_T=K)$, and $C$ is
    differentiable at $K>0$ if and only if $\Q(S_T=K)=0$.
  \item (\emph{Second differences and the law}) For $h>0$ and $K\in\R$,
    $\Delta_h^2C(K)=e^{-rT}\EQ[\Lambda_h(S_T-K)]\ge0$. As $h\downarrow0$:
    \begin{enumerate}[label=(\alph*)]
      \item $h\,\Delta_h^2C(K)\to e^{-rT}\Q(S_T=K)$ for every $K$;
      \item $\int_\R\phi(K)\,\Delta_h^2C(K)\,\dd K\to e^{-rT}\EQ[\phi(S_T)]$ for every bounded continuous
        $\phi:\R\to\R$: the measures $\Delta_h^2C(K)\,\dd K$ converge weakly to $e^{-rT}\mu$.
    \end{enumerate}
    Moreover $C''=e^{-rT}\mu$ in the sense of distributions:
    $\int_\R C(K)\,\phi''(K)\,\dd K=e^{-rT}\EQ[\phi(S_T)]$ for every $\phi\in C^2(\R)$ that vanishes
    outside a bounded interval.
  \item (\emph{Densities}) Let $I\subseteq(0,\infty)$ be an open interval. If $\mu$ has a density $q$
    on $I$, that is $\Q(S_T\in A)=\int_Aq(x)\,\dd x$ for Borel sets $A\subseteq I$, then $C$ is
    differentiable on $I$, $C'(b)-C'(a)=e^{-rT}\int_a^bq(x)\,\dd x$ for $a<b$ in $I$, and at every
    $K\in I$ where $q$ is continuous
    \[
      C''(K)=\lim_{h\downarrow0}\Delta_h^2C(K)=e^{-rT}q(K).
    \]
    Conversely, if $C$ is differentiable on $I$ and $C'$ is absolutely continuous on $I$ (for
    instance $C\in C^2(I)$), then $\mu$ has the density $e^{rT}C''$ on $I$.
  \item (\emph{Which functions are call prices}) A function $c:[0,\infty)\to\R$ is of the form
    $c(K)=e^{-rT}\int(x-K)^+\,\nu(\dd x)$ for a probability measure $\nu$ on $[0,\infty)$ with mean
    $F$ if and only if $c$ is convex, $c(0)=S_0e^{-yT}$, $\partial^+c(0)\ge-e^{-rT}$, and $c(K)\to0$ as
    $K\to\infty$. Then $\nu$ is unique: $\nu((K,\infty))=-e^{rT}\partial^+c(K)$ for $K\ge0$.
\end{enumerate}
\end{theorem}

\begin{proof}
(i) Let $K>0$ and $0<h\le K$. By \eqref{eq:sk-call-spread},
\begin{gather*}
  \frac{C(K+h)-C(K)}{h}=-e^{-rT}\EQ\Big[\frac{\min\big((S_T-K)^+,h\big)}h\Big],\\
  \frac{C(K)-C(K-h)}{h}=-e^{-rT}\EQ\Big[\frac{\min\big((S_T-K+h)^+,h\big)}h\Big].
\end{gather*}
Both integrands take values in $[0,1]$; as $h\downarrow0$ the first converges to $\1_{\{S_T>K\}}$ and
the second to $\1_{\{S_T\ge K\}}$, for every outcome. Dominated convergence gives the two
derivatives, and the first computation applies at $K=0$ as well.

(ii) The identity is \eqref{eq:sk-tent} at $x=S_T$, divided by $h^2$, and priced. (a) The variable
$h\Lambda_h(S_T-K)=(1-|S_T-K|/h)^+$ lies in $[0,1]$ and converges to $\1_{\{S_T=K\}}$; use dominated
convergence. (b) The integrand $|\phi(K)|\Lambda_h(S_T-K)$ is jointly measurable and its double
integral is at most $\sup|\phi|$, so by Fubini's theorem
\[
  \int_\R\phi(K)\,\Delta_h^2C(K)\,\dd K=e^{-rT}\EQ\Big[\int_\R\phi(S_T-u)\,\Lambda_h(u)\,\dd u\Big].
\]
For each outcome the inner integral is an average of $\phi$ over $[S_T-h,S_T+h]$, because $\Lambda_h$ is
a probability density on $[-h,h]$; it tends to $\phi(S_T)$ by continuity, and it is bounded by
$\sup|\phi|$. Dominated convergence concludes. For the distributional identity, let $\phi$ vanish
outside $[-M,M]$. Integrating by parts twice, with no boundary terms at $-M$,
$\int_\R(x-K)^+\phi''(K)\,\dd K=\int_{-M}^x(x-K)\phi''(K)\,\dd K=\phi(x)$ for every $x$. Since
$\int_\R|(x-K)^+\phi''(K)|\,\dd K\le2M(|x|+M)\sup|\phi''|$ and $\EQ|S_T|<\infty$, Fubini's theorem
lets us take $\EQ$ inside, which gives the identity.

(iii) $\mu$ has no atoms in $I$, so by (i) $C$ is differentiable on $I$ with
$C'(K)=-e^{-rT}\Q(S_T>K)$, and $C'(b)-C'(a)=e^{-rT}\Q(a<S_T\le b)=e^{-rT}\int_a^bq$. At a point $K$
where $q$ is continuous (in particular finite near $K$), the fundamental theorem of calculus gives
$C''(K)=e^{-rT}q(K)$; and for $h$ so small that $[K-h,K+h]\subset I$,
$\Delta_h^2C(K)=e^{-rT}\int q(x)\Lambda_h(x-K)\,\dd x$, an average of $q$ over $[K-h,K+h]$, which tends
to $e^{-rT}q(K)$. Conversely, if $C$ is differentiable on $I$, then $\mu$ has no atoms in $I$ by (i),
so $C'(K)=-e^{-rT}\Q(S_T>K)$ on $I$, and for $a<b$ in $I$, by absolute continuity,
$\Q(a<S_T\le b)=e^{rT}\big(C'(b)-C'(a)\big)=\int_a^be^{rT}C''(x)\,\dd x$; here $C''\ge0$ almost everywhere
because $C'$ is nondecreasing. The finite measures $\mu$ and $e^{rT}C''(x)\,\dd x$ on $I=(\alpha,\beta)$ agree on the
intervals $(a,b]\subset I$, hence, by continuity from below, on $(\alpha,b]$ for $b\in I$ and on $I$ itself.
Extended by zero outside $I$, they are finite measures on $\R$ with the same distribution function, so
they coincide (\textcite{durrett}, Theorem~1.1.4: a measure on $\R$ is determined by its Stieltjes
measure function).

(iv) \emph{Necessity.} $c$ is convex because each $K\mapsto(x-K)^+$ is, and integrating against $\nu$
preserves convexity; $c(0)=e^{-rT}\int x\,\nu(\dd x)=e^{-rT}F$; $\partial^+c(0)=-e^{-rT}\nu((0,\infty))\ge-e^{-rT}$ by (i),
whose proof only used the law; and $c(K)\to0$ by dominated convergence.

\emph{Sufficiency.} First, three facts about a convex $c$ on $[0,\infty)$ with $\partial^+c(0)$ finite.
It is continuous: in the interior because chord slopes are monotone, hence locally bounded, and at
$0$ because the right difference quotient has the finite limit $\partial^+c(0)$. Its right derivative
$g=\partial^+c$ exists everywhere, is nondecreasing, and is right-continuous: if $u_n\downarrow u$ and
$v>u_1$, then $g(u)\le g(u_n)\le(c(v)-c(u_n))/(v-u_n)$, and letting $n\to\infty$ and then $v\downarrow u$
gives $\lim_ng(u_n)=g(u)$. And $c(b)-c(a)=\int_a^bg(u)\,\dd u$: for $a=u_0<\dots<u_m=b$, convexity gives
$g(u_j)\le(c(u_{j+1})-c(u_j))/(u_{j+1}-u_j)\le g(u_{j+1})$, so $c(b)-c(a)$ lies between the lower and the
upper Riemann sums of the monotone, hence Riemann-integrable, function $g$.

Now $g\le0$, since $g(u_0)>0$ would force $c(u)\ge c(u_0)+g(u_0)(u-u_0)\to\infty$; and $g(u)\to0$ as
$u\to\infty$, since $g$ increases to some $L\le0$, and $L<0$ would force
$c(u)\le c(u_0)+L(u-u_0)\to-\infty$. So $G(u)=1+e^{rT}g(u)$ is nondecreasing, right-continuous, with
$G(0)\ge0$ by hypothesis and $G(u)\uparrow1$; set $G=0$ on $(-\infty,0)$. It is then the distribution
function of a probability $\nu$ on $[0,\infty)$ (\textcite{durrett}, Theorem~1.1.4), with
$\nu((u,\infty))=-e^{rT}g(u)$ for $u\ge0$. Writing $(x-K)^+=\int_K^\infty\1_{\{x>u\}}\,\dd u$ and
using Tonelli's theorem,
\[
  \int(x-K)^+\,\nu(\dd x)=\int_K^\infty\nu((u,\infty))\,\dd u=-e^{rT}\int_K^\infty g(u)\,\dd u
  =e^{rT}\Big(c(K)-\lim_{v\to\infty}c(v)\Big)=e^{rT}c(K).
\]
At $K=0$ this gives the mean $e^{rT}c(0)=F$. Uniqueness: by (i), every $\nu$ that produces $c$ has
$\nu((K,\infty))=-e^{rT}\partial^+c(K)$ for all $K\ge0$, so its distribution function is $G$, and a
measure is determined by its distribution function (\textcite{durrett}, Theorem~1.1.4).
\end{proof}

So the second strike derivative of the call is the discounted risk-neutral density when the
density exists and is continuous, and in general it is the discounted law itself, as a measure.
Prices at one maturity therefore identify the law of $S_T$ under $\Q$ (not the law under $\Prob$,
and not the law of the path), and part (iv) shows that \cref{prop:sk-strike-bounds} together
with $C(K)\to0$ is all that a single maturity can tell: every function with those properties is the
call price curve of exactly one law of $S_T$ with mean $F$. Whether that law is admissible for a given
$\Prob$ is a further condition: under (B) it must be equivalent to the $\Prob$-law of $S_T$
(\cref{rem:sk-finite}). In the standard model, (iii) gives
$C''(K)=G(0,S_0;T,K)$, the Green's function of \cref{prop:sm-green}.

Atoms show up in data through part (ii)(a): the butterfly $h^2\Delta_h^2C(K)$ centred at an
atom is asymptotic to $e^{-rT}\Q(S_T=K)\,h$, not $O(h^2)$, so $\Delta_h^2C$ has a spike of height close to
$e^{-rT}\Q(S_T=K)/h$ (\cref{fig:sk-strike}). The extreme case is a riskless stock, $\mu=\delta_F$:
then $C(K)=e^{-rT}(F-K)^+$, a ramp whose second derivative in the sense of (ii) is $e^{-rT}\delta_F$. In the
standard model the law of $S_T$ tends to $\delta_{S_0}$ as $T\downarrow0$ (\cref{prop:sm-green}), and the call
curve to the ramp $(S_0-K)^+$.

\begin{figure}[ht]
  \centering
  \includegraphics[width=\textwidth]{ch03-strike}
  \caption{Three laws of $S_T$ with the same mean $F=99.005$ (test case, $T=1$): lognormal with
    $\sigma=20\%$; the SSVI smile of \cref{sec:sk-smile}; and a takeover model in which $S_T=120$ with
    probability $p=0.2$ and is otherwise lognormal ($\sigma=20\%$, mean $93.756$). Left: $\Q(S_T>K)$
    computed as $-e^{rT}$ times the right slope of the call price, which jumps by $p$ at the atom.
    Right: $e^{rT}\Delta_h^2C(K)$ with $h=1$ on the strikes $40,41,\dots,180$, a density estimate for the
    first two laws and a spike of height $p/h$ plus the continuous part at $K=120$ for the third.}
  \label{fig:sk-strike}
\end{figure}

\demo{03-strike-arbitrage.ipynb}{move an atom of the law of $S_T$ and the strike step $h$, and
watch the kink in the call curve, the jump of the digital and the spike of the second
difference.}

\paragraph{From quotes to densities.} Quotes carry errors: rounding to the tick, the bid--ask
spread, stale prices. Second differences divide these errors by $h^2$, so the strike spacing has to
balance them against the bias of a coarse grid.

\begin{proposition}[Densities from noisy quotes]\label{prop:sk-fd-density}
Let $h>0$, let $C$ be four times continuously differentiable on $[K-h,K+h]$ with $|C''''|\le M_4$
there, and let $\tilde C$ be quotes with $|\tilde C-C|\le\varepsilon$ at $K-h$, $K$ and $K+h$. Then
\begin{equation}\label{eq:sk-fd-bound}
  \big|\Delta_h^2\tilde C(K)-C''(K)\big|\le\frac{M_4h^2}{12}+\frac{4\varepsilon}{h^2}.
\end{equation}
For $M_4,\varepsilon>0$ the right side is smallest at $h_*=(48\varepsilon/M_4)^{1/4}$, where it equals
$2\sqrt{M_4\varepsilon/3}$.
\end{proposition}

\begin{proof}
Taylor's formula with integral remainder at $K\pm h$, added, gives
$\Delta_h^2C(K)-C''(K)=h^{-2}\int_0^h\frac{(h-u)^3}{6}\big(C''''(K+u)+C''''(K-u)\big)\,\dd u$, whose absolute
value is at most $h^{-2}\cdot2M_4h^4/24=M_4h^2/12$. The errors contribute
$|e(K-h)-2e(K)+e(K+h)|/h^2\le4\varepsilon/h^2$, with $e=\tilde C-C$. The function $ah^2+b/h^2$ with
$a,b>0$ is smallest at $h^4=b/a$, where it equals $2\sqrt{ab}$.
\end{proof}

The bias is of order $h^2$ and the noise of order $\varepsilon/h^2$, so the best error is of order
$\sqrt\varepsilon$: halving the price errors improves the density by a factor $\sqrt2$ only. As an
example, take the SSVI prices of \cref{fig:sk-smile} at $T=1$, rounded to the cent
($\varepsilon=0.005$). On $[50,150]$, $\max|C''''|=2.763\times10^{-4}$, so $h_*=5.43$ and the bound is
$1.36\times10^{-3}$, against a peak value $0.0272$ of $C''$. On the strikes $60,61,\dots,140$ the largest
error of $\Delta_1^2\tilde C$ is $0.0169$, more than half the peak; on $60,65,\dots,140$ that of
$\Delta_5^2\tilde C$ is $0.0010$ (\cref{fig:sk-smile}, right). In practice one first fits an
arbitrage-free smile to the quotes and differentiates the fit (\cref{sec:sk-smile}); higher-order
stencils change the exponents (\exref{ex:sk-five-point}).

\section{The smile}\label{sec:sk-smile}

Let $C^{\mathrm{BS}}(K,\sigma)=S_0e^{-yT}N(d_+)-Ke^{-rT}N(d_-)$ be the Black--Scholes call
(\cref{thm:bsf-formula}), with $d_\pm=\big(\log(S_0/K)+(r-y\pm\sigma^2/2)T\big)/(\sigma\sqrt T)$, and let
$\mathcal V=S_0e^{-yT}N'(d_+)\sqrt T=Ke^{-rT}N'(d_-)\sqrt T$ be its vega (\cref{prop:bsf-greeks}). Under (B),
if $\Q(S_T<K)>0$ and $\Q(S_T>K)>0$, the bounds of \cref{prop:sk-strike-bounds}(i) are strict (if
$F\ge K$, the gap below is $P(K)>0$; if $F<K$, it is $C(K)>0$; the gap above is
$e^{-rT}\EQ[\min(S_T,K)]>0$), so by \cref{prop:bsf-implied-vol} there is exactly one $\sigma(K)>0$ with
$C(K)=C^{\mathrm{BS}}(K,\sigma(K))$: the implied volatility of \cref{def:bsf-implied-vol}. By parity the put
has the same one. The map $K\mapsto\sigma(K)$ is the \emph{smile} at maturity $T$. It is flat in the
Black--Scholes model, and it is not flat in the market.

\begin{remark}[The shape of equity-index smiles]\label{rem:sk-smile-data}
Since the crash of October 1987, implied volatilities of S\&P~500 index options decrease with the
strike over most of the quoted range: out-of-the-money puts are expensive relative to the lognormal
model \citep{rubinstein-1994}. The smile is therefore usually called a \emph{skew}; at short
maturities it often turns up again for far out-of-the-money calls. The skew is steepest at short
maturities and flattens as the maturity grows: \textcite{gatheral-volatility-surface}, Chapter~3,
shows the SPX surface of 15 September 2005 (Figure~3.2) and its at-the-money skew
$\partial\sigma/\partial\log K$ against the time to expiry (Figure~3.4). A traders' rule of thumb, used in
the course, scales the at-the-money skew like $T^{-1/2}$, so that smiles of different maturities nearly
collapse when plotted against $\log(K/S_0)/(\sigma_1\sqrt T)$, with $\sigma_1$ the one-year at-the-money
volatility. Over time the level of the surface moves with market stress more than its shape does.
The risk-neutral density that
\cref{thm:sk-breeden-litzenberger} extracts from a skew has a fatter left tail and a thinner right tail
than the lognormal one (\cref{fig:sk-smile,fig:sk-binary}), consistent with the heavy tails, the
gain/loss asymmetry and the negative correlation between returns and volatility (the leverage effect)
documented for index returns by \textcite{cont-stylized}, Section~2 (see also the SPX returns in
\textcite{gatheral-volatility-surface}, Figures~1.1--1.3), and with the premium that investors pay for
crash protection. Since quotes are proprietary, this chapter uses synthetic smiles of the same shape.
\end{remark}

Every strike derivative of $C(K)=C^{\mathrm{BS}}(K,\sigma(K))$ has a direct term, at fixed volatility, and
terms in the derivatives of the smile. For digitals this changes prices substantially.

\begin{proposition}[Digitals under a smile]\label{prop:sk-binary-smile}
Let $I\subseteq(0,\infty)$ be an open interval, $\sigma:I\to(0,\infty)$, and $C(K)=C^{\mathrm{BS}}(K,\sigma(K))$
for $K\in I$; below, $d_\pm$, $\mathcal V$ and $\sigma$ are evaluated at $\sigma=\sigma(K)$.
\begin{enumerate}[label=(\roman*)]
  \item If $\sigma$ is differentiable at $K\in I$, then so is $C$, and
    \begin{equation}\label{eq:sk-skew-term}
      -C'(K)=e^{-rT}N(d_-)-\mathcal V\,\sigma'(K).
    \end{equation}
    If moreover the prices come from a pricing measure (B), then $\Q(S_T=K)=0$, and the cash-or-nothing
    call paying $\1_{\{S_T>K\}}$ and the asset-or-nothing call paying $S_T\1_{\{S_T>K\}}$ cost
    \begin{gather*}
      D(K)=e^{-rT}\Q(S_T>K)=-C'(K),\\
      A(K)=e^{-rT}\EQ\big[S_T\1_{\{S_T>K\}}\big]=C(K)+K\,D(K)=S_0e^{-yT}N(d_+)-K\mathcal V\sigma'(K).
    \end{gather*}
  \item If $\sigma$ is differentiable on $I$ and twice differentiable at $K$, then $C''(K)$ exists and
    \begin{equation}\label{eq:sk-smile-density}
      C''(K)=\frac{e^{-rT}N'(d_-)}{K\sigma\sqrt T}\Big[1+2K\sqrt T\,d_+\sigma'+K^2T\,d_+d_-\,\sigma'^2+K^2T\sigma\,\sigma''\Big].
    \end{equation}
    Under (B), if $\sigma$ is twice continuously differentiable on $I$, the law of $S_T$ has the
    continuous density $e^{rT}C''$ on $I$.
\end{enumerate}
\end{proposition}

\begin{proof}
(i) $C^{\mathrm{BS}}$ is continuously differentiable in $(K,\sigma)$, with $\partial_KC^{\mathrm{BS}}=-e^{-rT}N(d_-)$ (the
terms in $\partial_Kd_\pm$ cancel because $S_0e^{-yT}N'(d_+)=Ke^{-rT}N'(d_-)$) and $\partial_\sigma C^{\mathrm{BS}}=\mathcal V$;
the chain rule gives \eqref{eq:sk-skew-term}. Under (B), \cref{thm:sk-breeden-litzenberger}(i) gives
$\Q(S_T=K)=0$ and $D(K)=-\partial^+C(K)=-C'(K)$. The identity
$(S_T-K)^+=S_T\1_{\{S_T>K\}}-K\1_{\{S_T>K\}}$ and the law of one price give $A=C+KD$, and
$C^{\mathrm{BS}}+Ke^{-rT}N(d_-)=S_0e^{-yT}N(d_+)$.

(ii) Differentiate $C'=\partial_KC^{\mathrm{BS}}+\partial_\sigma C^{\mathrm{BS}}\sigma'$ once more:
$C''=\partial_{KK}C^{\mathrm{BS}}+2\,\partial_{K\sigma}C^{\mathrm{BS}}\sigma'+\partial_{\sigma\sigma}C^{\mathrm{BS}}\sigma'^2+\mathcal V\sigma''$.
From $\partial_Kd_-=-1/(K\sigma\sqrt T)$ and $\partial_\sigma d_-=-d_+/\sigma$,
\[
  \partial_{KK}C^{\mathrm{BS}}=\frac{e^{-rT}N'(d_-)}{K\sigma\sqrt T},\qquad
  \partial_{K\sigma}C^{\mathrm{BS}}=\frac{e^{-rT}N'(d_-)\,d_+}{\sigma},\qquad
  \partial_{\sigma\sigma}C^{\mathrm{BS}}=\frac{\mathcal V\,d_+d_-}{\sigma},
\]
the last because $\mathcal V=Ke^{-rT}N'(d_-)\sqrt T$ and $\partial_\sigma N'(d_-)=N'(d_-)\,d_-d_+/\sigma$.
Factoring out $e^{-rT}N'(d_-)/(K\sigma\sqrt T)$ gives \eqref{eq:sk-smile-density}. Under (B) with
$\sigma\in C^2(I)$, also $C\in C^2(I)$, and \cref{thm:sk-breeden-litzenberger}(iii) applies.
\end{proof}

The \emph{skew term} $-\mathcal V\sigma'(K)$ in \eqref{eq:sk-skew-term} is the error of pricing a digital with
the Black--Scholes formula at the implied volatility of its own strike. For an equity skew,
$\sigma'<0$, so digital calls cost more and digital puts, $e^{-rT}-D(K)$, cost less than that formula
says. Second-order terms enter the density: \eqref{eq:sk-smile-density} contains the strike gamma,
the cross derivative $\partial_{K\sigma}C^{\mathrm{BS}}$ (counted twice) and the volga. With $\sigma'=\sigma''=0$
the bracket equals $1$, and \eqref{eq:sk-smile-density} is the discounted lognormal density, the Green's
function of \cref{prop:sm-green}.

\begin{example}[The course's smile]\label{exa:sk-course-digital}
The course models the one-year index smile near the money by
$\sigma(K)=\sigma_{\mathrm{atm}}+\xi\log(S_0/K)$, with $S_0=3900$, $T=1$, $r=3\%$, $y=1.3\%$,
$\sigma_{\mathrm{atm}}=25\%$ and $\xi=28\%$. At $K=S_0$, $\sigma'(K)=-\xi/K$, so the skew term is
$\xi\,e^{-yT}N'(d_+)\sqrt T$ and
\[
  D(S_0)=e^{-rT}N(d_-)+\xi e^{-yT}N'(d_+)\sqrt T=0.463167+0.108227=0.571394 .
\]
The skew term is $19\%$ of the price; leaving out the factor $e^{-yT}$, a common slip, gives $0.572810$.
The asset-or-nothing call costs $A(S_0)=413.04+3900\times0.571394=2641.47$, against $2219.39$ with a flat
$25\%$ volatility. \Cref{fig:sk-binary} shows the digitals and the densities across strikes: the smile
moves risk-neutral mass from the right tail to the left tail and shifts the mode to the right of the
forward.
\end{example}

\begin{figure}[ht]
  \centering
  \includegraphics[width=\textwidth]{ch03-binary}
  \caption{The course's smile $\sigma(K)=25\%+28\%\log(S_0/K)$ ($S_0=3900$, $T=1$, $r=3\%$, $y=1.3\%$)
    against a flat $25\%$ volatility. Left: digital call prices; the smile price $-C'(K)$ of
    \eqref{eq:sk-skew-term}, the Black--Scholes digital with a flat $25\%$, and the Black--Scholes
    digital $e^{-rT}N(d_-)$ at $\sigma(K)$, which omits the skew term. Right: the discounted density
    $C''(K)$ from \eqref{eq:sk-smile-density} and the discounted lognormal density with $\sigma=25\%$, both
    multiplied by $10^4$.}
  \label{fig:sk-binary}
\end{figure}

\paragraph{Testing a smile for arbitrage.} A smile $\sigma:(0,\infty)\to(0,\infty)$ is arbitrage-free at one
maturity when $C(K)=C^{\mathrm{BS}}(K,\sigma(K))$, with $C(0)=S_0e^{-yT}$, satisfies the conditions of
\cref{thm:sk-breeden-litzenberger}(iv). The conditions at $K=0$ hold automatically: Black--Scholes
prices satisfy $S_0e^{-yT}-Ke^{-rT}<C^{\mathrm{BS}}<S_0e^{-yT}$, so $C(K)\to C(0)$ and, once $C$ is convex,
$\partial^+C(0)\ge-e^{-rT}$. For a twice differentiable smile, what remains is that the bracket in
\eqref{eq:sk-smile-density} be nonnegative and that $C(K)\to0$ as $K\to\infty$. Written in the log-moneyness
$k=\log(K/F)$ and the total implied variance $w(k)=\sigma(K)^2T$, the bracket is the function
\[
  g(k)=\Big(1-\frac{k\,w'(k)}{2w(k)}\Big)^2-\frac{w'(k)^2}{4}\Big(\frac1{w(k)}+\frac14\Big)+\frac{w''(k)}2
\]
(\exref{ex:sk-durrleman}), and a slice is free of butterfly arbitrage if and only if $g\ge0$ and
$d_+\to-\infty$ as $k\to\infty$ (\textcite{gatheral-jacquier-2014}, eq.~(2.1) and Lemma~2.2).
\textcite{roper-arbitrage-free}, Theorems~2.9 and~2.15, shows for zero rates and dividends and for
twice-differentiable positive smiles that these conditions on every slice, together with total variance nondecreasing in $T$ and vanishing at $T=0$,
are sufficient and necessary for a surface free of static arbitrage; for constant $r$ and $y$ apply them to
the calls on $S_T/F_T$, $F_T=S_0e^{(r-y)T}$, whose prices are $C(K,T)/(S_0e^{-yT})$ at the strikes $K/F_T$. Arbitrage also bounds the
slope of a smile (\exref{ex:sk-skew-bounds}).

The course's smile fails the test far from the money. Its bracket is positive for
$84.33<K<9524.18$ and negative for $K<84.33$, that is below $2.2\%$ of the spot, where $\sigma=132\%$:
butterflies there have negative prices. Above $9524.18=S_0e^{\sigma_{\mathrm{atm}}/\xi}$ the formula gives
$\sigma\le0$. The total variance $(\sigma_{\mathrm{atm}}-\xi k')^2T$, $k'=\log(K/S_0)$, grows like $k'^2$, while under any
pricing measure $\limsup w(k)/|k|\le2$ in each wing: \textcite{lee-moment-formula}'s moment formula
(his Theorems~3.2 and~3.4) identifies the limsup as an explicit function, with values in $[0,2]$, of the
highest finite moment of $S_T$ (right wing) or of $1/S_T$ (left wing); see also
\textcite{gatheral-volatility-surface}, Chapter~7, ``Extreme strikes''. A near-the-money fit is not a smile on all of $(0,\infty)$, and digitals, densities and
Chapter~4's replication integrals reach far into the wings.

\paragraph{A synthetic surface.} The SSVI surface of \textcite{gatheral-jacquier-2014} has the right
shape and provably no arbitrage. With $\theta_T=\sigma_{\mathrm{atm}}^2T$ the at-the-money-forward total
variance, $\rho\in(-1,1)$ and a function $\phi>0$, it sets
\begin{equation}\label{eq:sk-ssvi}
  w(k,T)=\frac{\theta_T}2\Big(1+\rho\,\phi(\theta_T)k+\sqrt{\big(\phi(\theta_T)k+\rho\big)^2+1-\rho^2}\Big),\qquad
  \phi(\theta)=\frac{\eta}{\theta^\gamma(1+\theta)^{1-\gamma}}
\end{equation}
(their Definition~4.1 and eq.~(4.5)). We take the test case with $\sigma_{\mathrm{atm}}=20\%$, $\rho=-0.7$,
$\eta=1$, $\gamma=\frac12$. By their Remark~4.4, $\eta(1+|\rho|)\le2$ (here $1.7$) suffices for (4.5) to be free of
static arbitrage; directly, their Theorem~4.2 excludes butterfly arbitrage if
$\theta\phi(\theta)(1+|\rho|)<4$ and $\theta\phi(\theta)^2(1+|\rho|)\le4$ for all $\theta>0$; here
$\theta\phi(\theta)=\eta\sqrt{\theta/(1+\theta)}<1$ and $\theta\phi(\theta)^2=\eta^2/(1+\theta)<1$, so both left
sides are below $1.7$. Their Theorem~4.1 excludes calendar arbitrage if $\theta_T$ is nondecreasing and
$0\le\partial_\theta(\theta\phi)\le\phi\,(1+\sqrt{1-\rho^2})/\rho^2$; here $\partial_\theta(\theta\phi)/\phi=\frac12/(1+\theta)\le\frac12$,
and the right bound is $3.50$. The at-the-money skew is
$\partial_k\sigma=\rho\,\phi(\theta_T)\sqrt{\theta_T}/(2\sqrt T)\approx\rho\eta/(2\sqrt T)$ for small $\theta_T$: it scales
like $T^{-1/2}$, as in \cref{rem:sk-smile-data} ($-0.697$, $-0.343$ and $-0.238$ at $T=0.25$, $1$ and $2$).
At $T=1$ and $K=100$, $\sigma(K)=19.657\%$, the call costs $7.23033$ (against $7.36429$ at a flat $20\%$),
and the digital costs $0.436329+0.133830=0.570159$, against $0.436000$ at a flat $20\%$; the density of
$S_T$ at $100$ is $0.021125$, against $0.019724$ for the lognormal law.

\begin{figure}[ht]
  \centering
  \includegraphics[width=\textwidth]{ch03-smile}
  \caption{Left: implied volatilities of the SSVI surface \eqref{eq:sk-ssvi} ($\sigma_{\mathrm{atm}}=20\%$,
    $\rho=-0.7$, $\eta=1$, $\gamma=\frac12$; test case) against $\log(K/F_T)$ for $T=0.25$, $1$, $2$. Right: at $T=1$,
    the density $e^{rT}C''(K)$ of $S_T$ from \eqref{eq:sk-smile-density}, the lognormal density with
    $\sigma=20\%$, and the estimates $e^{rT}\Delta_h^2\tilde C(K)$ from SSVI prices rounded to the cent,
    with $h=1$ (strikes $60,61,\dots,140$) and $h=5$ (strikes $60,65,\dots,140$).}
  \label{fig:sk-smile}
\end{figure}

\demo{03-strike-arbitrage.ipynb}{change the SSVI parameters and the maturity, and compare the smile,
the digital with and without its skew term, the density and Durrleman's function $g$; then find
where the course's smile stops being arbitrage-free.}

\paragraph{Digitals from call spreads.} A digital is rarely listed, but call spreads are, and they
bound it. Let $h>0$ and let the strike grid be $h\N=\{0,h,2h,\dots\}$; a \emph{grid portfolio} holds bonds,
the prepaid forward and calls with strikes in $h\N$.

\begin{proposition}[Super- and sub-replication of a digital]\label{prop:sk-call-spreads}
Let $K\in h\N$ with $K\ge h$, and $\delta_K(x)=\1_{\{x>K\}}$.
\begin{enumerate}[label=(\roman*)]
  \item Among the payoffs $f(x)$ of grid portfolios with $f\ge\delta_K$ on $[0,\infty)$, the smallest
    one pointwise is $f^+(x)=\big((x-K+h)^+-(x-K)^+\big)/h$; among those with $f\le\delta_K$ on $[0,\infty)$,
    the largest is $f^-(x)=\big((x-K)^+-(x-K-h)^+\big)/h$.
  \item Under (P), the cheapest super-replicating grid portfolio costs
    $\bar D_h(K)=(C(K-h)-C(K))/h$ and the dearest sub-replicating one costs
    $\underline D_h(K)=(C(K)-C(K+h))/h$. If the digital is traded at the price $D(K)$ and (P) holds in the
    market that includes it, then $\underline D_h(K)\le D(K)\le\bar D_h(K)$.
  \item Under (B), $\bar D_h(K)-\underline D_h(K)=h\,\Delta_h^2C(K)$, which tends to $e^{-rT}\Q(S_T=K)$ as
    $h\downarrow0$; if the law of $S_T$ has a density bounded by $\bar q$ on $[K-h,K+h]$, the gap is at most
    $e^{-rT}\bar qh$.
\end{enumerate}
\end{proposition}

\begin{proof}
(i) The payoff $f$ of a grid portfolio is continuous and affine on each interval $[jh,(j+1)h]$. If
$f\ge\delta_K$, then $f\ge1$ on $(K,\infty)$, hence $f(K)\ge1$ by continuity, and $f(K-h)\ge0$; as $f$ is
affine on $[K-h,K]$, $f\ge f^+$ there, and $f\ge0=f^+$ on $[0,K-h]$ and $f\ge1=f^+$ on $[K,\infty)$. If
$f\le\delta_K$, then $f(K)\le0$ and $f(K+h)\le1$, so $f\le f^-$ on $[K,K+h]$, and $f\le0=f^-$ on $[0,K]$ and
$f\le1=f^-$ on $[K+h,\infty)$. Both $f^\pm$ are payoffs of call spreads on the grid. (ii) By (P), prices
are monotone in payoffs, so the extremes of (i) are the cheapest and the dearest; they also bound
the digital. (iii) The difference of the two spreads is $h$ times the butterfly \eqref{eq:sk-tent}
divided by $h^2$; apply \cref{thm:sk-breeden-litzenberger}(ii)(a). With a density $q\le\bar q$ on
$[K-h,K+h]$, $h\,\Delta_h^2C(K)=e^{-rT}\int q(x)\,h\Lambda_h(x-K)\,\dd x\le e^{-rT}\bar q\,h$, because
$\Lambda_h$ has integral $1$.
\end{proof}

The gap is first order in $h$, and it does not close at an atom. For the SSVI smile at $K=100$, the
bounds are $[0.5138,\,0.6184]$ with $h=5$ and $[0.5595,\,0.5804]$ with $h=1$, around $D(100)=0.5702$;
\exref{ex:sk-digital-grid} does the computation for the course's smile on its listed strike grids. A
dealer who sells the digital at $\bar D_h(K)$ and buys the spread carries no risk at all, whatever the
law of $S_T$; the price of this safety is the gap.

\begin{incode}
The module is \coderef{code/src/mfe230d/smile.py}. There, \texttt{strike\_checks} tests quotes against \eqref{eq:sk-checks};
\texttt{second\_difference}, \texttt{fd\_density}, \texttt{fd\_density\_error\_bound} and \texttt{fd\_optimal\_step}
implement \cref{prop:sk-fd-density}; \texttt{smile\_digital}, \texttt{smile\_asset\_digital} and
\texttt{smile\_density} implement \cref{prop:sk-binary-smile}; \texttt{course\_smile}, \texttt{ssvi\_smile},
\texttt{ssvi\_conditions} and \texttt{durrleman\_g} are the two synthetic smiles and their tests;
\texttt{call\_spread\_bounds} is \cref{prop:sk-call-spreads}; \texttt{atom\_call} prices calls on a law with an
atom. Black--Scholes prices and implied volatilities come from \coderef{code/src/mfe230d/bs.py}, and
\coderef{code/tests/test_smile.py} checks every number quoted in this chapter.
\end{incode}

\furtherreading
\textcite{breeden-litzenberger} is the original source of the density formula, and
\textcite{merton1973} of the strike relations. \textcite{gatheral-volatility-surface} is the
practitioner's account of the volatility surface; its Chapter~8 discusses digital options under a skew.
\textcite{rubinstein-1994} reads risk-neutral laws from index options and builds trees that fit them.
\textcite{davis-hobson}, \textcite{carr-madan-2005} and \textcite{roper-arbitrage-free} characterize prices and
surfaces free of static arbitrage, \textcite{lee-moment-formula} bounds the wings of a smile, and
\textcite{gatheral-svi-2004} and \textcite{gatheral-jacquier-2014} give the SVI and SSVI parameterizations.
